\section{Анализ предметной области и обоснование темы исследования}

\subsection{Тенденции развития многоядерных и многопроцессорных архитектур}

На протяжении нескольких десятилетий основным способом повышения производительности процессоров являлось увеличение тактовой частоты. Данная стратегия позволяла разработчикам программного обеспечения практически бесплатно получать прирост производительности при переходе на новое поколение оборудования. Однако в начале 2000-х годов этот подход столкнулся с фундаментальными физическими ограничениями: дальнейший рост тактовой частоты приводил к недопустимому увеличению тепловыделения и энергопотребления~\cite{Sutter2005}.

Ответом индустрии на данное ограничение стал переход к многоядерным архитектурам. Вместо наращивания производительности одного ядра производители начали размещать на одном кристалле несколько вычислительных ядер, каждое из которых работает на относительно невысокой тактовой частоте. Согласно данным компании Intel, уже к 2006 году большинство новых потребительских процессоров выпускалось с двумя и более ядрами, а к 2023 году флагманские серверные процессоры содержат сотни вычислительных ядер~\cite{IntelXeon2023}.

Данная тенденция нашла отражение в известном высказывании Херба Саттера, ставшем крылатой фразой в сообществе разработчиков: «Бесплатный обед закончился»~\cite{Sutter2005}. Суть этого утверждения состоит в том, что простое обновление оборудования более не даёт автоматического прироста производительности для однопоточных программ. Вместо этого разработчики вынуждены явно проектировать свои приложения с расчётом на параллельное выполнение, чтобы задействовать все доступные вычислительные ресурсы.

Параллельно с развитием многоядерных процессоров активно развивались и многопроцессорные системы, в которых несколько физических процессоров объединяются в единую вычислительную платформу. Такие системы получили широкое распространение в области высокопроизводительных вычислений (HPC), центрах обработки данных и серверных инфраструктурах крупных интернет-компаний. Архитектура NUMA (Non-Uniform Memory Access), характерная для многопроцессорных систем, вносит дополнительную сложность в разработку параллельных программ: время доступа к памяти зависит от того, на каком узле она физически расположена~\cite{Lameter2013}.

Таким образом, к середине 2010-х годов параллелизм превратился из узкоспециализированной темы, интересной лишь разработчикам суперкомпьютеров, в повседневную реальность для широкого круга инженеров-программистов. Веб-серверы, игровые движки, системы машинного обучения, базы данных и пользовательские приложения --- все они вынуждены эффективно использовать параллелизм для достижения приемлемой производительности~\cite{Herlihy2020}.

Рост числа ядер порождает принципиально новые требования к механизмам управления параллельным выполнением задач. Простые решения, приемлемые при двух-четырёх ядрах, демонстрируют неудовлетворительное масштабирование при 32, 64 и более ядрах. Именно поэтому исследование эффективных механизмов управления потоками и распределения задач между ними приобретает особую актуальность в современных условиях.

\subsection{Проблемы наивного подхода к управлению потоками}

Наиболее очевидным и простым подходом к параллельному выполнению задач является создание нового потока для каждой поступающей задачи и его уничтожение по завершении работы. Данный подход интуитивно понятен, прост в реализации и нередко используется в учебных примерах и прототипах. Однако при применении в производственных системах он обнаруживает ряд серьёзных недостатков.

Высокие накладные расходы на создание и уничтожение потоков. Создание потока операционной системой является дорогостоящей операцией. В операционной системе Linux создание потока через вызов \texttt{clone()} занимает порядка 5--15 мкс~\cite{Drepper2003}. В операционной системе Windows аналогичная операция через \texttt{CreateThread()} имеет сопоставимую стоимость. При этом каждый поток требует выделения стека (как правило, от 64 КБ до 8 МБ по умолчанию), инициализации структур данных ядра и регистрации в планировщике операционной системы. Для высоконагруженных систем, обрабатывающих тысячи и десятки тысяч задач в секунду, совокупные накладные расходы на управление жизненным циклом потоков могут составлять значительную долю общего времени выполнения~\cite{Butenhof1997}.

Неконтролируемый рост числа потоков. В условиях пиковой нагрузки наивный подход приводит к лавинообразному росту числа одновременно существующих потоков. Поскольку каждый поток потребляет память (прежде всего под стек) и ресурсы ядра операционной системы, бесконтрольное создание потоков способно привести к исчерпанию памяти или достижению системного лимита на количество потоков. В Linux данный лимит определяется параметром ядра \texttt{/proc/sys/kernel/threads-max} и по умолчанию составляет порядка нескольких тысяч потоков. Превышение этого лимита влечёт отказ в создании новых потоков и, как следствие, отказ в обслуживании запросов~\cite{Love2013}.

Избыточное переключение контекста. Планировщик операционной системы распределяет процессорное время между всеми активными потоками посредством вытесняющей многозадачности. Переключение контекста --- сохранение состояния текущего потока и восстановление состояния следующего --- является затратной операцией, требующей нескольких микросекунд и приводящей к инвалидации кэшей процессора. Когда число активных потоков значительно превышает количество ядер, система тратит существенную долю времени на переключение контекстов вместо полезной работы~\cite{Tanenbaum2014}. Данный эффект получил название thrashing и является одной из ключевых проблем наивного подхода.

Отсутствие контроля над использованием ресурсов. Наивный подход не предоставляет механизмов ограничения количества одновременно выполняемых задач. В результате в периоды пиковой нагрузки система оказывается перегруженной, что негативно сказывается на латентности обработки всех запросов, включая приоритетные. Концепция backpressure --- механизма обратного давления, при котором производитель задач замедляется при переполнении системы --- принципиально несовместима с наивным подходом~\cite{Reactive2014}.

Неэффективное использование кэша процессора. Потоки, созданные для выполнения коротких задач, как правило, не имеют устойчивой привязки к конкретным ядрам процессора. Планировщик ОС может перемещать поток между ядрами при каждом пробуждении, что приводит к систематическим промахам кэша и снижению эффективности работы с памятью. В современных процессорах разница в латентности между кэшем L1 и основной памятью достигает 200--300 раз, поэтому влияние данного фактора на производительность весьма существенно~\cite{Drepper2007}.

Перечисленные недостатки наивного подхода делают его непригодным для применения в производственных системах с высокими требованиями к производительности и предсказуемости. Это обстоятельство обусловило широкое распространение альтернативной концепции --- пула потоков, при которой фиксированный набор заранее созданных потоков многократно используется для выполнения поступающих задач. Однако и в рамках концепции пула потоков остаётся ряд нерешённых проблем, связанных с балансировкой нагрузки между потоками, что будет рассмотрено в последующих разделах.

